\documentclass[]{article}
\usepackage[portuguese]{babel}
\usepackage{graphicx}
\usepackage{tikz}
\usetikzlibrary{automata, positioning, arrows}
\usepackage[americanvoltages] {circuitikz}
%opening

\begin{document}

\textbf{Resolução de exercícios do livro}

\textbf{Data:} 21/09/2026

\textbf{Nome:} Leandro Júnior Cândido de Oliveira
\newline
\hrule

\begin{center}
	\textbf{Questões}
\end{center}
\begin{figure}[htp]
	\centering

\begin{tikzpicture}[node distance=1.5cm, shorten >=1mm, shorten <=0.5mm]
	% Estados
	\node[state, initial, accepting, thick, fill=gray!50, initial text=] (q1) {\textbf{q1}};
	\node[state, thick, fill=gray!50]            (q2) [below left=of q1]{\textbf{q2}};
	\node[state, thick, fill=gray!50]            (q3) [below right=of q1]{\textbf{q3}};
	
	%Transições
	
	\path[->, thick] (q1) edge node [above left] {b} (q2)
					(q1) edge [bend right=20] node [left] {$\epsilon$} (q3)
					(q3) edge [bend right=20] node [above right] {a} (q1)
					(q2) edge [loop, out=130, in=180, looseness=4] node [above left]{a} (q2)
					(q2) edge node [below] {a, b} (q3);
\end{tikzpicture}
\caption{Autômato finito AFN.}
\label{fig1:AutAFN}
\vspace{5pt}
\footnotesize Fonte: Leandro (2026).
\end{figure}

\end{document}
